\documentclass[12pt]{article}
\usepackage[a4paper,margin=1in]{geometry}
\usepackage[T1]{fontenc}
\usepackage[utf8]{inputenc}

\title{SRS V2 - RF1, RF2V2 y RF3V2}
\author{}
\date{}

\begin{document}
\maketitle

\section{Introduccion}
Este SRS describe los requisitos funcionales RF1, RF2V2 y RF3V2 para el sistema TekMess. El objetivo es asegurar la gestion de empleados, el acceso al sistema y el reporteo analitico.

\section{RF1: Gestion de Empleados}
\subsection{Proposito}
Permitir la gestion de empleados con operaciones de crear, editar y eliminar, manteniendo datos consistentes.

\subsection{Entradas}
\begin{itemize}
  \item Cedula
  \item Nombres
  \item Rol
  \item Correo institucional
\end{itemize}

\subsection{Salidas}
\begin{itemize}
  \item Confirmacion de registro, edicion o baja
  \item Credenciales generadas cuando aplica
\end{itemize}

\subsection{Reglas de negocio}
\begin{itemize}
  \item Cedula valida (10 digitos)
  \item Rol RBAC permitido
  \item Correo institucional valido
\end{itemize}

\section{RF2V2: Autenticacion simple}
\subsection{Proposito}
Permitir el acceso al sistema con usuario y contrasena. Incluye cambio y recuperacion de contrasena sin verificacion externa.

\subsection{Casos cubiertos}
\begin{itemize}
  \item Iniciar sesion
  \item Cambiar contrasena
  \item Recuperar contrasena
  \item Validar credenciales
\end{itemize}

\subsection{Entradas}
\begin{itemize}
  \item Usuario
  \item Contrasena
  \item Correo (solo para recuperacion)
\end{itemize}

\subsection{Salidas}
\begin{itemize}
  \item Acceso concedido o denegado
  \item Confirmacion de cambio o recuperacion
\end{itemize}

\section{RF3V2: Reporte}
\subsection{Proposito}
Generar y consultar reportes con informacion operativa y analitica basada en RF1 y RF2V2. El acceso a reportes depende del rol.

\subsection{Contenido del reporte}
\begin{itemize}
  \item Cantidad de empleados gestionados
  \item Cantidad de empleados editados
  \item Cantidad de accesos fallidos
\end{itemize}

\subsection{Origen de datos}
\begin{itemize}
  \item Gestion de empleados (RF1)
  \item Autenticacion y accesos (RF2V2)
\end{itemize}

\subsection{Entradas}
\begin{itemize}
  \item Rango de fechas o periodo de consulta
\end{itemize}

\subsection{Salidas}
\begin{itemize}
  \item Reporte analitico generado
  \item Historial de reportes consultados
\end{itemize}

\section{Trazabilidad}
\begin{itemize}
  \item RF3V2 depende de RF1 para datos de empleados.
  \item RF3V2 depende de RF2V2 para datos de acceso.
  \item RF2V2 controla el acceso segun rol.
\end{itemize}

\end{document}
